\documentclass{beamer}
\usetheme{Rochester}
\usepackage[french]{babel}
\usepackage{tikz,lstautogobble,listings, minted}
\usepackage[T1]{fontenc}
\usepackage[utf8]{inputenc}

\usetikzlibrary{trees,shapes,arrows.meta,positioning, automata, calc}
\setbeamertemplate{navigation symbols}{}
\setbeamertemplate{footline}[text line]{%
\parbox{\linewidth}{\insertshortauthor\hfill\insertframenumber / \inserttotalframenumber}}


\setminted[ocaml]{	                               
	linenos=true, numbersep=4pt, xleftmargin=20pt,      %--numéro de ligne
	breaklines=true                                    %--découpage des longues lignes
}
\setminted[python]{	                               
	linenos=true, numbersep=4pt, xleftmargin=20pt,      %--numéro de ligne
	breaklines=true                                    %--découpage des longues lignes
}
\setminted[promela]{	                               
	linenos=true, numbersep=4pt, xleftmargin=20pt,      %--numéro de ligne
	breaklines=true                                    %--découpage des longues lignes
}

\theoremstyle{definition}
\newtheorem{mydef}[theorem]{Définition}

\title{Méthode des tableaux: Applications dans différentes logiques}
\author{GIL Dorian - Candidat n°20220}
\subtitle{TIPE 25/26 - Cycles et Boucles}
\date{}

\begin{document}

\begin{frame}
\titlepage
\end{frame}

\begin{frame}{Contexte}
    \textbf{Objectif:} Explorer l'intérêt dans l'industrie d'un algorithme abstrait: la Méthode des tableaux.
    \pause
    \vspace{10px}
    
    En particulier:
    \begin{itemize}
        \item Problème SAT
        \item Vérification de Modèle
    \end{itemize}
\end{frame}

\begin{frame}{Définitions}
    \begin{mydef}[Logique Propositionelle]
        On la définit inductivement par
        \begin{itemize}
            \item $V, F$ sont des formules
            \item les variables booléennes sont des formules
            \item Si $\alpha$ est une formule, $\lnot\alpha$ (NON alpha) est une formule
            \item Si $\alpha, \beta$ est une formule, $\alpha\land\beta$ (ET), $\alpha\lor\beta$ (OU) sont des formules.
        \end{itemize}
    \end{mydef}
    \begin{mydef}[Litteraux]
        Variable ou la négation d'une variable.
    \end{mydef}

\end{frame}

\begin{frame}{Comment résoudre SAT ?}
    \begin{center}
        \textbf{La Méthode des tableaux (Logique Propositionelle)}

        \textit{Algorithme de résolution du problème SAT}
    \end{center}
    \pause
    \begin{center}
        \textbf{Règle d'inférence:}    
    \end{center}
    \begin{center}
    \begin{tikzpicture}[scale=0.75]
    \node {$\phi\lor\psi$}
        child {node {$\phi$}}
        child {node {$\psi$}};
    \end{tikzpicture}
    \begin{tikzpicture}[scale=0.75]
    \node {$\lnot(\phi\land\psi)$}
        child {node {$\lnot\phi$}}
        child {node {$\lnot\psi$}};
    \end{tikzpicture}
    \begin{tikzpicture}[scale=0.75]
    \node {$\lnot\lnot\phi$}
        child {node {$\phi$}};
    \end{tikzpicture}
    \begin{tikzpicture}[scale=0.75]
    \node {$\phi\land\psi$}
        child {node {$\phi, \psi$}};
    \end{tikzpicture}  
    \begin{tikzpicture}[scale=0.75]
    \node {$\lnot(\phi\lor\psi)$}
        child {node {$\lnot\phi, \lnot\psi$}};
    \end{tikzpicture}
    \end{center}
    \pause
    \textbf{Exemple avec $a\land \lnot a$ :}

    \begin{tikzpicture}[scale=0.75]
    \node {$a\land\lnot a$}
        child {node {$a, \lnot a$}};
    \end{tikzpicture} 
\end{frame}

\begin{frame}{Étude introductive}
    \begin{mydef}[Forme Alternée]
        Soit $n\in\mathbb{N}^*$, et $(a_k)_{k\in [|1,n|]}$ des litteraux, on dit que $\varphi$ est de forme alternée ssi
        $\varphi = a_1\land(a_2\lor(a_3\land(\dots(a_n))))$
    \end{mydef}
    \pause
    \begin{center}
        \textbf{Objectif de l'étude introductive ? (Optimisable ?, Interêt)} 
    \end{center}
    \pause
    \begin{center}
        \textbf{Une re-écriture:} 
    \end{center}
    \begin{columns}[t]
    \begin{column}{0.30\textwidth}
        \centering
        \begin{tikzpicture}[level distance=8mm]
        \node {$a_1$}
            child {node {$a_2$}}
            child {node {$a_3$}
            child {node {$a_4$}}
            child {node {$a_5$}}};
        \end{tikzpicture}
    \end{column}
   
    \begin{column}{0.10\textwidth}
    \centering
    \begin{tikzpicture}[x=1mm,y=1mm]
      \draw[<-, thick, black] (-10,0) .. controls +(4,0) and +(-4,-4) .. (4,-4);
    \end{tikzpicture}
  \end{column}

    \begin{column}{0.40\textwidth}
        \centering
        \begin{tikzpicture}[level distance=8mm]
        \node {$a_1$}
            child {node {$a_1 , a_2\lor(a_3\land(a_4\lor a_5))$}
                child {node {$a_2$}}
                child {node {$a_3\land(a_4\lor a_5)$}
                    child {node {$a_3 , a_4\lor a_5$}
                        child {node {$a_4$}}
                        child {node {$a_5$}}
                        }
                }
            };
        \end{tikzpicture}
    \end{column}
    \end{columns}
\end{frame}

\begin{frame}[fragile]{Résolution}
    \textbf{Proposition d'algorithme ($\mathcal{O}(n)$, correction prouvé):}
    \begin{enumerate}
        \item Fin de l'arbre: Renvoyer Vrai
        \item Feuille gauche: SI Pas Contradiction, Renvoyer Vrai
        \item Feuille droite: SI Contradiction, Renvoyer Faux SINON Appel Recursif
    \end{enumerate}
    \begin{tikzpicture}[level distance=8mm]
        \node {$a$}
            child {node {$\lnot a$}}
            child {node {$b$}
            child {node {$c$}}
            child {node {$d$}}};
    \end{tikzpicture}
    \pause

    \vspace{10px}
    Pour 100 Formes alternées, temps d'execution:
    \begin{itemize}
        \item \textbf{Alternée:} 0.000493s
        \item \textbf{Quine (avec conversion en CNF):} 0.025874s
        \item \textbf{Quine (sans conversion en CNF):} 0.018691s
    \end{itemize} 

\end{frame}
 
\begin{frame}{Conclusion intermédiaire}
    \textbf{Quelques Remarques}
    \begin{itemize}[<+->]
        \item Conversion IMPOSSIBLE (Exemple: $a_1\lor a_2$)
        \item Interêt SAT dans l'industrie
        \item Interêt de l'algorithme pour SAT ?
        \item Réelle application de la méthode ?
    \end{itemize}
\end{frame}

\begin{frame}{Contexte}
    \begin{center}
        \textbf{La Méthode des tableaux (Logique Temporelle Linéaire)}
        \textit{Algorithme utilisé en vérification de modèle}
    \end{center}
    \begin{mydef}[Verification de modèle (Model Checking)]
    Technique automatique qui étant donné un modèle à nombre d'état fini d'un système et des propriétés formelles,
    vérifie systématiquement si ces propriétés restent vrai pour ce modèle.
    \end{mydef}
    \begin{mydef}[Logique Linéaire Temporelle (LTL)]
    Logique propositionnelle avec dépendance discrète en temps
    \end{mydef}
    \pause
    \textbf{Pourquoi la LTL ?} (Raisonnement sur des circuits intégrés...)
\end{frame}

\begin{frame}{Principe de l'algorithme}
    \textbf{Hypothèse:}
    \begin{itemize}
        \item On admet l'existence de algorithme
        \item Il utilise la méthode des tableaux
    \end{itemize}
    \pause
    \vspace{10px}
    \textbf{Idée de l'algorithme}
    \begin{itemize}
        \item Des conditions (Formule logique en LTL)
        \item Un modèle (Système de Transition)
        \item Renvoie un retour sur la compatibilité des deux
    \end{itemize}
    \vspace{10px}
    Appliquons l'algorithme sur le système suivant:
    
    \textbf{Compteur Binaire Modulo 4: } Tout les quatres cycles, il renvoie 1 (vrai) sinon 0 (faux).
\end{frame}

\begin{frame}{Système de Transitions}
    \begin{mydef}[Système de transition]
        C'est un automate qui prend en paramètre
        \begin{itemize}
            \item $S$ un ensemble d'états, $I$ les états initiaux
            \item $Act$ un ensemble d'actions
            \item $\rightarrow$ une relation de transition ($S\times Act \rightarrow S$)
            \item $AP$ un ensemble de variable
            \item $L$ une fonction qui a un état associe une partie de $AP$
        \end{itemize} 
    \end{mydef}

    \centering
    \begin{tikzpicture}[>=Stealth]
    % nodes at explicit coordinates
    \node[draw,rounded corners,minimum width=1.2cm,minimum height=0.7cm] (s00) at (0,1.5) {00};
    \node[draw,rounded corners,minimum width=1.2cm,minimum height=0.7cm] (s01) at (3,1.5) {01};
    \node[draw,rounded corners,minimum width=1.2cm,minimum height=0.7cm] (s10) at (3,0)   {10};
    \node[draw,rounded corners,minimum width=1.2cm,minimum height=0.7cm] (s11) at (0,0)   {11};

    % small annotations (use math mode for underscores/comma)
    \node at (-1,1.4) {$\{r\}$};
    \node at (4,1.4)  {$\{\varphi_2\}$};
    \node at (4,0) {$\{\varphi_1\}$};
    \node at (-1.3,0) {$\{\varphi_1, \varphi_2\}$};

    % edges (square)
    \draw[->] (s00) -- (s01);
    \draw[->] (s01) -- (s10);
    \draw[->] (s10) -- (s11);
    \draw[->] (s11) -- (s00);

    

    % self-loop on s00
    \draw[->] (-0.5,2) -- (s00.north);
    \end{tikzpicture}
\end{frame}

\begin{frame}{Contraintes}
    \textbf{Propriétés}
    \begin{itemize}
        \item L'ampoule $2$ change toujours de valeur
        \item L'ampoule $1$ change ssi $\varphi_2$
        \item $r$ ssi les deux ampoules sont éteintes
    \end{itemize}
    \pause
    \textbf{Traduction}
    \begin{itemize}
        \item $G(\varphi_2 \Leftrightarrow X(\lnot\varphi_2))$
        \item $G(\varphi_2\Leftrightarrow(\varphi_1 \Leftrightarrow X(\lnot\varphi_1)))$
        \item $G(r \Leftrightarrow \lnot(\varphi_1\lor\varphi_2))$
    \end{itemize}
    \pause
    \textbf{Remarque: } 
    \begin{itemize}
        \item $G$ assure que la propriété reste vraie en tout temps
        \item $X\alpha$ = Est ce que $\alpha$ est vraie au temps prochain ? ($X$ pour neXt)
    \end{itemize}
\end{frame}

\begin{frame}{Comparaison}
    Source: Principles of Model Checking de C.Baier et JP.Katoen
    \begin{itemize}
        \item Même système de transition
        \item $G(r\Leftrightarrow\lnot\varphi_1\land\lnot\varphi_2)$ similaire
        \item $G(\varphi_1 \implies(Xr \lor XXr))$
        \item $G(r\lor Xr\lor XXr \lor XXXr)$
        \item $G(r\implies(X\lnot r \land XX\lnot r \land XXX \lnot r))$
    \end{itemize}
    
    \pause
    \textbf{Remarques:}
    \begin{itemize}
        \item A priori, cycle modulo 4 pas assuré par mes contraintes
        \item Leur contrainte laisse libre l'ordre d'allumage des lumières
        \item Ce n'est pas forcement une preuve de correction
    \end{itemize}
\end{frame}

\begin{frame}{Exemple d'execution}
    \textbf{Outil Test:} SPIN
    
    Compatibilité avec $G(r \Leftrightarrow \lnot(\varphi_1\lor\varphi_2))$ ?
    \vspace{10px}

    \textbf{Resultat (extrait):}

    + Partial Order Reduction\\
    full statespace search for:\\
    never claim + (p3)\\
    assertion violations + (if within scope of claim)\\
    acceptance cycles + (fairness disabled)\\
    invalid end states - (disabled by never claim)\\
    \vspace{5px}

    17 states, stored\\
    17 states, matched\\
    0 atomic steps\\
    hash conflicts: 0 (resolved)
\end{frame}

\begin{frame}{Conclusion}
    Rappel: On voulait déterminer l'utilité de la méthode des tableaux dans différentes logiques.
    \begin{itemize}
        \item Étude de la résolution de SAT et familiarisation avec la méthode 
        \item Application et discussion à propos des algorithmes utilisant cette méthode
    \end{itemize}
\end{frame}

\inputminted[]{ocaml}{Code/Alternee_Tableau.ml}
\inputminted[]{python}{Code/PropGenerator.py}
\inputminted[]{ocaml}{Code/quine.ml}
\inputminted[]{promela}{Code/modulo4.pml}

\end{document}